\section{Bruno Díaz Garnica}
La elaboración de este proyecto fue un muy buen cierre a la clase de Robótica, ya que nos permitió integrar todos los conocimientos aprendidos a lo largo del semestre, pero esta vez enfocados en nuestro propio robot personalizado, lo cual involucró descubrir nuevos tipos de retos. Una de las principales dificultades fue el aprendizaje de nuevos programas nunca antes vistos en la carrera, como ROS y Gazebo, así como el uso más extensivo de MATLAB. 

Personalmente, uno de los problemas más destacados durante la elaboración del proyecto fue desarrollar la Tabla de Denavit Hartenberg del robot, ya que ésta era la base para la simulación y prueba de los diferentes códigos de cinemática; por lo tanto, si la tabla era incorrecta significaba que ningún resultado era correcto. Otro problema destacado, fue que en la prueba de cinemática inversa cometimos un error al inicio, por lo que el tiempo que tardaba en correr el código era muy extenso, complicando el proceso de checar trayectorias. 

Resolver estos errores fue un gran aprendizaje para la resolución general de problemas, ya que nos motivó a investigar las causas de las problematicas y trabajar en equipo para idear una solución adecuada. 